\documentclass[10pt]{article}
\usepackage[T1]{fontenc}
\usepackage[utf8]{inputenc}
\usepackage[french]{babel}
\usepackage{amsmath,amssymb,mathrsfs,mathpazo}
\usepackage{geometry}
\geometry{a4paper,margin=2cm}
\usepackage{graphicx,array,tabularx,tabularray,colortbl,multirow,multicol,booktabs,enumitem,pifont,ulem,fancybox,diagbox,cancel}
\usepackage{pstricks,pstricks-add,pst-plot,pst-tree,pst-node,pst-text,pst-func,pst-math,pst-bspline,pst-3dplot,pst-eucl,pst-solides3d}
\let\pscalculate\undefined
\usepackage{pst-bezier}
\usepackage{tikz,pgfplots,tkz-tab}
\usetikzlibrary{arrows,arrows.meta,patterns,calc,positioning,trees,decorations.markings,shapes.geometric,shapes.misc,shapes.arrows,intersections,decorations.pathmorphing,decorations.pathreplacing,backgrounds,shadows,babel}
\pgfplotsset{compat=1.18}
\usepgfplotslibrary{fillbetween}
\usepackage[np]{numprint}
\usepackage{listings,makecell,pas-tableur,twoopt}
\let\barre\undefined
\usepackage[pstricks]{tablvar}
\definecolor{darkgreen}{RGB}{0,100,0}
\providecommand{\red}{\color{red}}
\providecommand{\blue}{\color{blue}}
\providecommand{\green}{\color{green}}
\providecommand{\black}{\color{black}}
\providecommand{\white}{\color{white}}
\providecommand{\gray}{\color{gray}}
\providecommand{\cyan}{\color{cyan}}
\providecommand{\magenta}{\color{magenta}}
\providecommand{\yellow}{\color{yellow}}
\providecommand{\violet}{\color{violet}}
\providecommand{\orange}{\color{orange}}
\providecommand{\euro}{\texteuro}
\providecommand{\EUR}{\texteuro}
\providecommand{\ieme}{\textsuperscript{e}}
\providecommand{\degres}{\ensuremath{^\circ}}
\providecommand{\degree}{\ensuremath{^\circ}}
\providecommand{\vv}[1]{\overrightarrow{#1}}
\providecommand{\Coord}[1]{#1}
\providecommand{\decosix}{}
\providecommand{\starredbullet}{\ensuremath{\bullet}}
\pagestyle{empty}
\setlength{\parindent}{0pt}
\usepackage{adjustbox}\DeclareUnicodeCharacter{2236}{:}\DeclareUnicodeCharacter{2190}{\ensuremath{\leftarrow}}
\providecommand{\R}{\mathbb{R}}
\providecommand{\N}{\mathbb{N}}
\providecommand{\Z}{\mathbb{Z}}
\providecommand{\Q}{\mathbb{Q}}
\providecommand{\D}{\mathbb{D}}
\providecommand{\C}{\mathbb{C}}
\providecommand{\e}{{\mathrm{e}}}
\providecommand{\vect}[1]{\overrightarrow{#1}}
\providecommand{\vectt}[1]{\overrightarrow{\mathrm{#1}}}
\providecommand{\vvt}[1]{\overrightarrow{\mathrm{#1}}}
\providecommand{\np}[2][]{{#2~#1}}
\providecommand{\vv}[1]{{\overrightarrow{#1}}}
\renewcommand{\d}{{\mathrm{d}}}
\providecommand{\degres}{{^\circ}}
\providecommand{\degre}{{^\circ}}
\providecommand{\blue}{\color{blue}}
\providecommand{\red}{\color{red}}
\providecommand{\enskip}{\quad}
\providecommand{\llbracket}{\lbrack\!\lbrack}
\providecommand{\rrbracket}{\rbrack\!\rbrack}
\providecommand{\euro}{{\text{€}}}
\providecommand{\no}{{\text{n°}}}
\providecommand{\interval}[2]{{[#1;#2]}}
\providecommand{\intervFO}[2]{{[#1;#2[}}
\providecommand{\dotfill}{\ldots}

\providecommand{\ds}{\displaystyle}
\providecommand{\textding}[1]{#1}
\providecommand{\Oij}{\(\mathrm{O};\vec{\imath},\vec{\jmath}\)}
\providecommand{\Oijk}{\(\mathrm{O};\vec{\imath},\vec{\jmath},\vec{k}\)}
\providecommand{\Ouv}{\(\mathrm{O};\vec{u},\vec{v}\)}
\providecommand{\intervFF}[2]{\ensuremath{[#1;#2]}}
\renewcommand{\euro}{\ifmmode\text{\texteuro}\else\texteuro\fi}
\DeclareUnicodeCharacter{00B0}{\ifmmode^{\circ}\else\textdegree\fi}
\DeclareUnicodeCharacter{20AC}{\ifmmode\text{\texteuro}\else\texteuro\fi}
% Énoncés : Éducation nationale. Transcription et corrigés : APMEP, Denis Vergès et contributeurs. Découpage et classement : CoopMaths, Sébastien Lozano, Éric Elter et contributeurs. Adaptation de la présentation : Maths Explorer.
% CC BY-SA : https://coopmaths.fr/www/about/
% https://coopmaths.fr/alea/static/bac/2026/tex/bac_2026_06_sujet2_antilles_4.tex
\begin{document}
\noindent\textit{APMEP / CoopMaths. CC BY-SA. Présentation adaptée par Maths Explorer.}\par\medskip
On considère la fonction $f_{k}$ définie sur $] 0 ;+\infty[$ par :
	\[f_{k}(x)=1-x+\dfrac{k(1+\ln x)}{x}\]
	où $k$ est un réel strictement positif et ln désigne la fonction logarithme népérien.

	\medskip

	On désigne par $\mathcal{C}_{k}$ la représentation graphique de la fonction $f_{k}$ dans un repère orthonormé.

	\medskip

	On admet que la fonction $f_{k}$ est deux fois dérivable sur $] 0 ;+\infty[$.

	\medskip

	On note $f_{k}'$ sa dérivée première et $f_{k}''$ sa dérivée seconde.

	\subsection*{Partie A : étude d'un premier cas particulier}
		Dans cette partie, on se place dans le cas $k=2$.

		La courbe $\mathcal{C}_{2}$ est représentée sur le graphique ci-dessous.

		On précise que :

		\begin{itemize}[leftmargin=10mm]
			\item Le point A est un point d'inflexion de la courbe $\mathcal{C}_{2}$;
			\item $T_{2}$ est la tangente à la courbe $\mathcal{C}_{2}$ au point A.

			\begin{center}
			\begin{tikzpicture}[baseline={(current bounding box.center)}]
				\begin{axis}[/pgf/number format/.cd, use comma,
					x={22mm}, y={22mm},
					xmin = -0.3, xmax = 3.6,
					ymin = -1.6, ymax = 2.4,
					xtick={0,1,...,4},
					ytick ={-2,...,2},
					xticklabels={0,1},
					yticklabels={,,0,1},
					minor tick num = 1,
					grid=both, axis lines =center]
					\addplot [line width=1.pt,color=teal,smooth,samples=100,domain= 0.1:3.6,dashed]{-1.369*(x-1.649)+1.171};
					\addplot [line width=1.2pt,color=red,smooth,samples=100,domain= 0.1:3.6]{1-x+2*(1+ln(x))/x};
					\node[below left] at (axis cs: 0,0) {{0}};
					\node[red] at(axis cs: 0.47,-1.25) {$ \mathcal{C}_2$};
					\node[teal] at(axis cs: 1,2.25) {$T_2$};
					\draw[shift={(axis cs: 1.649,1.171)}] (2pt,-2pt)--(-2pt,2pt) (-2pt,-2pt)--(2pt,2pt) node[above right]{A} ;
				\end{axis}
			\end{tikzpicture}
			\end{center}
		\end{itemize}

		\begin{enumerate}
			\item On considère une primitive $F_{2}$ de la fonction $f_{2}$ sur $] 0 ;+\infty[$.

			Avec la précision permise par le graphique, donner les variations de la fonction $F_{2}$ sur l'intervalle $]0 ; 3]$.

			\item Parmi les trois courbes ci-dessous, indiquer celle qui représente $f_{2}'$ et celle qui représente $f_{2}''$.

			Justifier.
		\end{enumerate}

		\begin{tabular}{ccc}
			\begin{tikzpicture}[baseline={(current bounding box.center)}]
				\begin{axis}[/pgf/number format/.cd, use comma,
					x={11mm}, y={11mm},
					xmin = -0.3, xmax = 3.6,
					ymin = -1.6, ymax = 2.4,
					xtick={0,1,...,4},
					ytick ={-2,...,2},
					xticklabels={0,1},
					yticklabels={,,0,1},
					minor tick num = 1,
					grid=both, axis lines =center]
					\addplot [line width=1.2pt,color=red,smooth,samples=100,domain= 0.5:3.6]{-1-2*ln(x)/(x^2)};
					\node[below left] at (axis cs: 0,0) {{0}};
				\end{axis}
			\end{tikzpicture}&
			\begin{tikzpicture}[baseline={(current bounding box.center)}]
				\begin{axis}[/pgf/number format/.cd, use comma,
					x={11mm}, y={11mm},
					xmin = -0.3, xmax = 3.6,
					ymin = -1.6, ymax = 2.4,
					xtick={0,1,...,4},
					ytick ={-2,...,2},
					xticklabels={0,1},
					yticklabels={,,0,1},
					minor tick num = 1,
					grid=both, axis lines =center]
					\addplot [line width=1.2pt,color=red,smooth,samples=100,domain= 0.9:3.6]{(-2+4*ln(x))/(x^3)};
					\node[below left] at (axis cs: 0,0) {{0}};
				\end{axis}
			\end{tikzpicture}&
			\begin{tikzpicture}[baseline={(current bounding box.center)}]
				\begin{axis}[/pgf/number format/.cd, use comma,
					x={11mm}, y={11mm},
					xmin = -0.3, xmax = 3.6,
					ymin = -1.6, ymax = 2.4,
					xtick={0,1,...,4},
					ytick ={-2,...,2},
					xticklabels={0,1},
					yticklabels={,,0,1},
					minor tick num = 1,
					grid=both, axis lines =center]
					\draw[line width=1.2pt,color=red,smooth]
					(axis cs: 0.5,-1.6) .. controls (axis cs: 1,-0.85) and (axis cs: 1.5,-0.4) .. (axis cs: 2.5,0)
					.. controls (axis cs: 3.5,0.4) and (axis cs: 3.6,0.35) .. (axis cs: 3.6,0.35);
					\node[below left] at (axis cs: 0,0) {{0}};
				\end{axis}
			\end{tikzpicture}\\
			Courbe 1 & Courbe 2 & Courbe 3\\
		\end{tabular}

	\subsection*{Partie B : étude d'un second cas particulier}
		Dans cette partie on se place dans le cas $k=1$.

		\medskip

		On considère la fonction $f_{1}$ définie sur $] 0 ;+\infty[$ par :
		\[f_{1}(x)=1-x+\dfrac{1+\ln x}{x}\]

		\begin{enumerate}
			\item \begin{enumerate}
				\item Montrer que pour tout réel $x>0$ :
				\[f_{1}'(x)=-1-\dfrac{\ln x}{x^{2}}.\]

				\item Déterminer la limite de $f_{1}'(x)$ lorsque $x$ tend vers $+\infty$.
			\end{enumerate}

			\item On admet que :
			\begin{itemize}[leftmargin=1cm]
				\item l'axe des ordonnées est une asymptote de la courbe de $f_{1}'$;
				\item pour tout réel $x>0$, \quad $f_{1}''(x)=\dfrac{2 \ln x-1}{x^{3}}$.
			\end{itemize}

			On donne ci-dessous le tableau de variations de $f_{1}'$ sur l'intervalle $] 0 ;+\infty[$.

			\begin{center}
			\begin{tikzpicture}
				\tkzTabInit[lgt=1,espcl=4]
				{$x$/0.8, $f_{1}'$/2}
				{0,$\e^{\frac{1}{2}}$,$+\infty$}
				\tkzTabVar{D+/$+\infty$,-/$-1-\dfrac{1}{2\e}$,+/$-1$}
			\end{tikzpicture}
			\end{center}

			\begin{enumerate}
				\item Justifier les variations de la fonction $f_{1}'$ sur l'intervalle $]0 ;+\infty[$ et la valeur de l'extremum.

				\item Justifier la limite de la fonction $f_{1}'$ en 0.

				\item Montrer qu'il existe une unique solution $\alpha$ à l'équation \quad $f_{1}'(x)=0$ \quad sur $] 0 ;+\infty[$.

				En déduire que $f_{1}$ admet un maximum sur l'intervalle $] 0 ;+\infty[$.

				\item Recopier et compléter les pointillés du programme écrit en langage Python ci-dessous, afin qu'il renvoie une valeur approchée à $10^{-2}$ près de la solution $\alpha$ de l'équation \quad $f_{1}'(x)=0$ \quad sur $] 0 ;+\infty[$.

				\begin{center}
				\begin{tblr}{colspec={X[6cm, font=\ttfamily]},rowsep=0pt, vlines, hline{1,Z}={solid}}
					def d(x) :\\
					\quad return -1-ln(x)/x**2\\
					\\
					def alpha() :\\
					\quad u = 0.5\\
					\quad while d(u) ... 0 :\\
					\quad \quad u = ...\\
					\quad return u\\
				\end{tblr}
				\end{center}

				On rappelle que \texttt{x**2} désigne $x^{2}$.
			\end{enumerate}

			\end{enumerate}

	\subsection*{Partie C: $\boldsymbol{k}$ quelconque}
		On considère la fonction $h$ définie sur $] 0 ;+\infty[$ par \quad
		$h(x)=\dfrac{1+\ln x}{x}$,

		et la fonction $H$ définie sur $] 0 ;+\infty[$ par \quad
		$H(x)=\dfrac{1}{2}(1+\ln x)^{2}$.

		On admet que $H$ est dérivable sur $] 0 ;+\infty[$.

		\begin{enumerate}
			\item Montrer que la fonction $H$ est une primitive de $h$.
		\end{enumerate}

		On rappelle que $k$ est un réel strictement positif et que $f_{k}$ est la fonction définie sur $] 0 ;+\infty[$ par :
		\[f_{k}(x)=1-x+\dfrac{k(1+\ln x)}{x}\]

		\begin{enumerate}
			\setcounter{enumi}{1}
			\item Existe-t-il une valeur de $k$ pour laquelle la valeur moyenne de la fonction $f_{k}$ sur l'intervalle $[ 1 ; \e]$ est égale à 0 ?

			Justifier.
		\end{enumerate}

\end{document}
